\documentclass[a4paper, serif, 10pt]{moderncv}

%% moderncv
% color and style
\moderncvcolor{black}
\moderncvstyle{banking}

%% page layout/margins
\usepackage[top=2.5cm, bottom=2.5cm, left=1.5cm, right=1.5cm]{geometry}

\nopagenumbers{}

%% Babel config for Dutch language
% ref:
% - https://latex3.github.io/babel/guides/locale-dutch.html
\usepackage[dutch]{babel}

%% fontspec
\usepackage{fontspec}

\IfFontExistsTF{Linux Libertine}{
  \setmainfont[Ligatures={TeX, Common}, Numbers=OldStyle]{Linux Libertine}
}{}
\IfFontExistsTF{Linux Libertine O}{
  \setmainfont[Ligatures={TeX, Common}, Numbers=OldStyle]{Linux Libertine O}
}{}

%% CTeX
% CJK support, used to show CJK characters in the resume
%
% - heading=false: disable ctex heading
% - punct=kaiming: use kaiming punctuations styles for CJK
% - scheme=plain: use plain scheme, do not override \normalsize font size
% - space=auto: space settings for CJK characters
%
% CJK fonts are resolved per locale language, e.g. Noto Serif/Sans CJK JP
% for Japanese resumes, with fallback to the next region if not installed
%
% ref:
% - http://ctan.mirrorcatalogs.com/language/chinese/ctex/ctex.pdf
\usepackage[UTF8, heading=false, punct=kaiming, scheme=plain, space=auto]{ctex}

\IfFontExistsTF{Noto Serif CJK SC}{
  \setCJKmainfont{Noto Serif CJK SC}
}{}
\IfFontExistsTF{Noto Sans CJK SC}{
  \setCJKsansfont{Noto Sans CJK SC}
}{}

%% line spacing
\usepackage{setspace}
\setstretch{1.125}

%% URL styling - use normal text instead of monospace
\urlstyle{same}

%% Patch \httplink and \httpslink to handle full URLs
% The original moderncv macros always prepend http:// or https:// to URLs.
% This causes issues when the URL already has a protocol (e.g., https://example.com)
% resulting in malformed URLs like https://https://example.com.
% This patch checks if the URL already contains "://" and uses it directly if so.
% Note: We use \str_set:Nx (x-expansion) to expand macros like \@homepage first.
\ExplSyntaxOn
\str_new:N \g_yamlresume_url_str

\RenewDocumentCommand{\httpslink}{O{}m}{
  \str_set:Nx \g_yamlresume_url_str {#2}
  \str_if_in:NnTF \g_yamlresume_url_str {://}
    { \url{#2} }
    { \url{https://#2} }
}

\RenewDocumentCommand{\httplink}{O{}m}{
  \str_set:Nx \g_yamlresume_url_str {#2}
  \str_if_in:NnTF \g_yamlresume_url_str {://}
    { \url{#2} }
    { \url{http://#2} }
}
\ExplSyntaxOff

\name{Sanne de Vries}{}
\title{Senior HR-manager met focus op mens en cultuur}
\phone[mobile]{+31 6 1234 5678}
\email{sanne.devries@example.nl}
\homepage{https://www.sannedevries.nl}

\address{Prinsengracht 456, Amsterdam, Noord-Holland, Nederland, 1016 HK}{}{}

\extrainfo{{\small \faLinkedin}\ \href{https://www.linkedin.com/in/sannedevries}{@sannedevries} {} {} {} • {} {} {} 
{\small \faMedium}\ \href{https://medium.com/@sannedevries}{@sannedevries}}

\begin{document}

\maketitle

\section{Basis Informatie}

\cvline{}{\begin{itemize}
\item Senior HR-manager met meer dan tien jaar ervaring in strategisch HRM, talentmanagement en organisatieontwikkeling.
\item Bewezen trackrecord in het vertalen van bedrijfsdoelen naar personeelsagenda's met meetbare impact op verloop, betrokkenheid en productiviteit.
\item Sterk in stakeholdermanagement, arbeidsrecht en verandermanagement binnen complexe, internationale organisaties.
\item Datagedreven en mensgericht; coacht leidinggevenden en bouwt HR-teams die proactief bijdragen aan de bedrijfsresultaten.
\end{itemize}}
\section{Opleidingen}

\cventry{sep 2011–jun 2013}
        {Master, Arbeids- en Organisatiepsychologie, Score: 8,2}
        {Universiteit van Amsterdam}
        {\url{https://www.uva.nl}}
        {}
        {\begin{itemize}
\item Afgestudeerd op de relatie tussen leiderschapsstijl en betrokkenheid in hybride teams.
\item Verdieping in arbeids- en organisatiepsychologie, onderzoeksmethoden en HR-analytics.
\item Scriptie beoordeeld met een 8,7.
\end{itemize}
\textbf{Cursussen}: Arbeids- en organisatiepsychologie, Strategisch HRM, Onderzoeksmethoden, Verandermanagement, Arbeidsrecht}

\cventry{sep 2007–jun 2011}
        {Bachelor, Human Resource Management, Score: 8,5}
        {Hogeschool van Amsterdam}
        {\url{https://www.hva.nl}}
        {}
        {\begin{itemize}
\item Breed fundament in HRM, arbeidsrecht en organisatiekunde.
\item Stage bij een regionale zorgorganisatie op het gebied van werving en selectie.
\end{itemize}
\textbf{Cursussen}: Human Resource Management, Arbeidsrecht, Organisatiekunde, Communicatie, Personeelsplanning}
\section{Werkervaring}

\cventry{mrt 2021–Heden}
        {Senior HR-manager}
        {PeopleHub Nederland}
        {\url{https://www.peoplehub.nl}}
        {}
        {\begin{itemize}
\item Ontwikkelt en implementeert de HR-strategie voor een organisatie van 750 medewerkers.
\item Stuurt een team van vijf HR-professionals aan en realiseert een jaarlijkse daling van 18\% in verloop.
\item Adviseert directie en MT over organisatieontwikkeling, talentmanagement en arbeidsrechtelijke vraagstukken.
\item Trekt het HR-analytics-programma en vertaalt data naar concrete interventies voor behoud en ontwikkeling.
\end{itemize}
\textbf{Trefwoorden}: HR-strategie, Talentmanagement, HR-analytics, Leiderschapsontwikkeling, Arbeidsrecht}

\cventry{sep 2016–feb 2021}
        {HR-manager}
        {Van Dijk Retail Groep}
        {\url{https://www.vandijkretail.nl}}
        {}
        {\begin{itemize}
\item Verantwoordelijk voor HR-beleid, werving en selectie, onboarding en prestatiemanagement voor 12 vestigingen.
\item Begeleidde een reorganisatie met behoud van 95\% van kritische functies.
\item Introduceerde een medewerkersonderzoek met een respons van 89\% en koppelde resultaten aan verbeterplannen.
\item Onderhandelde met vakbonden over CAO-afspraken en sociale plannen.
\end{itemize}
\textbf{Trefwoorden}: Reorganisatie, Werving en selectie, Medezeggenschap, CAO, Prestatiemanagement}

\cventry{sep 2013–aug 2016}
        {HR-adviseur}
        {TalentWerk Advies}
        {\url{https://www.talentwerkadvies.nl}}
        {}
        {\begin{itemize}
\item Adviseerde operationele managers in logistiek en customer service over HR-vraagstukken.
\item Ontwikkelde een onboardingprogramma dat de time-to-productivity met drie weken verkortte.
\item Behandelde verzuim- en reintegratietrajecten in samenwerking met bedrijfsarts en leidinggevenden.
\item Voerde HR-analyses uit en rapporteerde maandelijks over verzuim, verloop en bezetting.
\end{itemize}
\textbf{Trefwoorden}: Onboarding, Verzuimbegeleiding, HR-administratie, Adviesvaardigheden}
\section{Talen}

\cvline{Nederlands}{Moedertaal of tweetalige vaardigheid \hfill \textbf{Trefwoorden}: Moedertaal, Zakelijk schrijven}
\cvline{Engels}{Volledige professionele vaardigheid \hfill \textbf{Trefwoorden}: TOEFL iBT 108, Zakelijke communicatie}
\cvline{Duits}{Beperkte werkvaardigheid \hfill \textbf{Trefwoorden}: Goethe-Zertifikat B1}
\section{Vaardigheden}

\cvline{Strategisch HRM}{Expert \hfill \textbf{Trefwoorden}: HR-strategie, Organisatieontwikkeling, Personeelsplanning}
\cvline{Talentmanagement}{Geavanceerd \hfill \textbf{Trefwoorden}: Werving en selectie, Onboarding, Successieplanning, Prestatiemanagement}
\cvline{Arbeidsrecht en compliance}{Geavanceerd \hfill \textbf{Trefwoorden}: CAO, Arbeidsomstandigheden, Medezeggenschap, Privacy (AVG)}
\cvline{Verandermanagement}{Geavanceerd \hfill \textbf{Trefwoorden}: Cultuurverandering, Stakeholdermanagement, Lean}
\cvline{HR-analytics}{Intermediair \hfill \textbf{Trefwoorden}: Power BI, Excel, HR-dashboards, Data-gedreven besluitvorming}
\cvline{Communicatie}{Expert \hfill \textbf{Trefwoorden}: Coaching, Onderhandelen, Presenteren, Adviesvaardigheden}
\section{Onderscheidingen}

\cventry{nov 2023}
        {Nederlandse Vereniging voor Personeelsmanagement}
        {HR Professional van het Jaar 2023}
        {}
        {}
        {\begin{itemize}
\item Prijswinnaar in de categorie strategisch HRM vanwege het programma Duurzame Inzetbaarheid.
\item Jury prees de datagedreven aanpak en meetbare impact op verloop en betrokkenheid.
\end{itemize}}
\section{Certificaten}

\cventry{mei 2020}
        {SHRM}
        {SHRM-SCP}
        {\url{https://www.shrm.org/certification}}
        {}
        {}

\cventry{okt 2019}
        {FNV}
        {Gecertificeerd Medezeggenschapsadviseur}
        {\url{https://www.fnv.nl}}
        {}
        {}

\cventry{mrt 2022}
        {Coursera}
        {HR Analytics Fundamentals}
        {\url{https://www.coursera.org}}
        {}
        {}
\section{Publicaties}

\cventry{apr 2022}
        {Wendbaar werken: van ambitie naar routines}
        {Kluwer HRM}
        {\url{https://www.kluwer.nl/hrm}}
        {}
        {\begin{itemize}
\item Praktijkboek over het vertalen van strategische HR-doelen naar dagelijkse routines.
\item Met casuïstiek uit retail, zorg en technologie.
\end{itemize}}
\section{Referenties}

\cventry{\emaillink[m.jansen@peoplehub.nl]{m.jansen@peoplehub.nl}}
        {Directeur HR bij PeopleHub Nederland}
        {Marieke Jansen}
        {+31 20 123 4567}
        {}
        {Sanne is een strategische en daadkrachtige HR-manager die complexe HR-vraagstukken vertaalt naar heldere plannen en meetbare resultaten. Zij combineert scherpe analyse met oprechte aandacht voor mensen en is een waardevolle sparringpartner voor elk managementteam.}
\section{Projecten}

\cventry{jan 2022–dec 2022}
        {Organisatiebreed programma om werkplezier, vitaliteit en blijvende inzetbaarheid van medewerkers te versterken.}
        {Programma Duurzame Inzetbaarheid}
        {\url{https://www.peoplehub.nl/projecten/duurzame-inzetbaarheid}}
        {}
        {\begin{itemize}
\item Ontwikkelde een meetbaar interventieportfolio voor 750 medewerkers.
\item Realiseerde een daling van het verzuim met 1,8 procentpunt en een stijging van de betrokkenheidsscore met 12 punten.
\item Werkte samen met bedrijfsarts, OR en lijnmanagement aan preventie en reintegratie.
\end{itemize}
\textbf{Trefwoorden}: Vitaliteit, Verzuimpreventie, Betrokkenheid, HR-analytics}
\section{Interesses}

\cvline{Vrijwilligerswerk}{Loopbaancoaching, Mentoring}
\cvline{Sport}{Hardlopen, Wielrennen, Yoga}
\cvline{Cultuur}{Literatuur, Musea, Toneel}
\section{Vrijwilligerswerk}

\cventry{sep 2018–Heden}
        {Loopbaanmentor}
        {Stichting LoopbaanLab}
        {\url{https://www.loopbaanlab.nl}}
        {}
        {\begin{itemize}
\item Begeleidt jaarlijks vijf professionals bij loopbaanvragen en re-integratie.
\item Verzorgt workshops over personal branding en solliciteren.
\item Draagt bij aan de ontwikkeling van laagdrempelige loopbaantools.
\end{itemize}}
    
\end{document}